\documentclass[11pt,a4paper]{article}
\usepackage[utf8]{inputenc}
\usepackage[T1]{fontenc}
\usepackage[spanish,es-noshorthands]{babel}
\usepackage[margin=2.5cm]{geometry}
\usepackage{booktabs}
\usepackage{enumitem}
\usepackage{hyperref}

\title{Resumen de la bitácora de diseño\\\large Proyecto \emph{canasta-scraper}}
\author{Analítica de la Web}
\date{20 de septiembre de 2026}

\begin{document}
\maketitle

\begin{abstract}
Este documento resume las decisiones registradas en \texttt{BITACORA.md}: qué se
decidió, con qué evidencia y qué quedó abierto. Los números citados se midieron
sobre los datos del 19 y 20 de septiembre de 2026 y deben recalcularse sobre el
panel acumulado antes de entrar al paper.
\end{abstract}

\section{Qué es el EAN}

El \textbf{EAN} (\emph{European Article Number}, hoy parte del estándar GS1) es el
código numérico de 13 dígitos (a veces 8 o 12) que aparece bajo el código de
barras de un producto envasado. Lo asigna el fabricante, no la tienda, y por eso
tiene dos propiedades que lo vuelven central para el proyecto:

\begin{itemize}[nosep]
  \item \textbf{Es único por producto físico.} Un mismo EAN en Metro, Wong, Plaza
        Vea, Vivanda o Tottus identifica exactamente la misma presentación
        (misma marca, mismo gramaje, mismo envase). Es la \emph{llave de
        emparejamiento} entre cadenas: los \texttt{productId} de cada tienda no
        son comparables entre sí, el EAN sí.
  \item \textbf{Es estático.} El código de barras no cambia con el tiempo, así
        que basta consultarlo una vez por SKU y guardarlo en caché
        (\texttt{data/ean\_tottus.json}), a diferencia del precio, que hay que
        capturar todos los días.
\end{itemize}

No todas las cadenas lo publican igual: Metro lo expone en el 100\,\% de sus
productos, Plaza Vea en el 58\,\%, y Tottus solo en la ficha de producto, no en el
listado. Los productos sin EAN (marca propia, granel, panadería) son justamente
la parte del catálogo que no se puede emparejar por ahora. Además, el EAN
\emph{regala etiquetas}: pares con el mismo EAN son ``mismo producto'' y pares con
distinto EAN son ``distinto producto'', lo que permite entrenar y evaluar un
modelo de emparejamiento sin anotar nada a mano.

\section{Mapa del pipeline}

Flujo: \texttt{config/retailers.yml} $\rightarrow$ \texttt{collect.py} $\rightarrow$
cliente (\texttt{vtex} $|$ \texttt{falabella\_catalyst}) $\rightarrow$
\texttt{normalize} $\rightarrow$ \texttt{storage}. Se guarda en dos capas: JSON crudo
comprimido y CSV normalizado. Las categorías de alimentos se filtran
\textbf{por nombre}, no por id.

Tres piezas donde vive la dificultad real:
\begin{enumerate}[nosep]
  \item \texttt{parse\_quantity()}: el gramaje se extrae del nombre; el modo de
        falla peligroso es el gramaje \emph{presente y equivocado}.
  \item \texttt{eans.py}: caché persistente y versionada; si se pierde, nunca converge.
  \item \texttt{storage.py}: el crudo existe para reprocesar cuando aparezca un bug.
\end{enumerate}

\textbf{Principio de fondo:} lo único irrecuperable es el precio de hoy. Todo lo
demás (EAN, gramaje, emparejamiento, precio con tarjeta) se repara después sobre
el crudo. De ahí la tolerancia a fallos: una cadena caída no tumba la corrida.

\section{Qué se descartó de la propuesta original}

\begin{table}[h]
\centering
\begin{tabular}{p{5cm}p{9.5cm}}
\toprule
\textbf{Elemento} & \textbf{Realidad} \\
\midrule
Makro & No existe: tienda online muerta. Sacar de la slide. \\
``4 súper (si alcanza: Vivanda)'' & Ya alcanzó. Son \textbf{5 cadenas}. \\
IPC-Web vs.\ IPC del INEI como núcleo & \textbf{No viable}: para diciembre habrá 2--3 observaciones mensuales del INEI. \\
Comparación con el INEI en general & Confundida por canal: el INEI releva mercados de abastos; el panel, solo supermercado. \\
``50--100 alimentos'' & Número inventado. Hay \textbf{917} productos con EAN en las 5 cadenas. \\
\bottomrule
\end{tabular}
\end{table}

Se conservan: el esquema de 22 columnas, el precio por kg, \texttt{on\_sale} vs.\
inflación, los tres antecedentes (Cavallo \& Rigobon 2016; Jaworski 2021; Bueno
2024) y la fórmula de la canasta.

Diagnóstico de fondo: ``un tablero de esta semana'' describe un \emph{producto},
no una pregunta de investigación. Un tablero no puede estar equivocado; un paper
necesita una afirmación falseable.

\section{Pregunta de investigación redefinida}

\begin{quote}
\textbf{¿El precio de un mismo producto físico (mismo EAN) en la web de los
supermercados de Lima se fija por cadena, por grupo corporativo o por
competencia?}
\end{quote}

Sub-preguntas: (1) dispersión, (2) estructura de propiedad y formato,
(3) rigidez y liderazgo de precios, (4) promoción vs.\ inflación.

\subsection*{Evidencia con un día de datos (2026-09-20, replicado el 19)}

\begin{table}[h]
\centering
\begin{tabular}{lrrr}
\toprule
\textbf{Par de cadenas} & \textbf{EAN comunes} & \textbf{Precio idéntico} & \textbf{Brecha mediana} \\
\midrule
Metro + Wong (Cencosud)       & 5\,444 & 67.7\,\% & 0.0\,\% \\
Wong + Vivanda                & 1\,987 & 46.5\,\% & 1.8\,\% \\
Plaza Vea + Tottus            & 3\,081 & 44.7\,\% & 1.4\,\% \\
Plaza Vea + Vivanda (SPSA)    & 4\,725 & 23.8\,\% & 2.4\,\% \\
Vivanda + Tottus              & 2\,466 & 14.4\,\% & 5.6\,\% \\
\bottomrule
\end{tabular}
\end{table}

Dispersión (brecha max/min): mediana 4.0\,\% con EAN en $\geq 2$ cadenas
(11\,019 EAN); 13.9\,\% en los 906 EAN presentes en las 5 cadenas (p90: 41.4\,\%).

\textbf{Dos hallazgos:}
\begin{itemize}[nosep]
  \item Los grupos corporativos no se comportan igual: Cencosud sirve el mismo
        precio en Metro y Wong 2 de cada 3 veces (backend de pricing compartido);
        SPSA no, porque Vivanda está posicionada premium.
  \item Plaza Vea coincide más con Tottus (44.7\,\%) que con su hermana Vivanda
        (23.8\,\%): la competencia disciplina más el precio que la propiedad común.
\end{itemize}

\textbf{Por qué es mejor:} desriesga (las sub-preguntas 1 y 2 ya están
contestadas); convierte ``solo mido supermercado'' en irrelevante; convierte
``Metro y Wong son la misma empresa'' en la variable independiente; es más
Analítica de la Web (dos storefronts sirviendo el mismo precio desde el mismo
backend; asimetría de publicación del EAN); la canasta sobrevive entera como
entregable de cierre.

\section{Dónde entra el modelo}

\subsection*{Dónde no}
Predecir el precio de mañana: el 99.5\,\% de los precios no cambia de un día a
otro, así que ``mañana igual que hoy'' acierta $\sim$99\,\% sin aprender nada.

\subsection*{Emparejamiento sin EAN: preprocesamiento, no modelamiento}
Cerca de la mitad del universo no es emparejable por EAN (13\,119 de 24\,188
productos con EAN en una sola cadena, más $\sim$10\,700 filas sin EAN). Usando
pares con EAN compartido como verdad de campo (Metro vs.\ Plaza Vea, 2\,294 pares):
nombre idéntico tras normalizar 3.5\,\%; Jaccard $\geq 0.8$ solo 24\,\% de
recall, con 0 falsos positivos. El texto es preciso pero ciego; el 76\,\% restante
requiere modelo (TF-IDF como baseline, Sentence-BERT siamés con pérdida
contrastiva sobre $\sim$31k pares, gradient boosting con gramaje y precio como
features). \textbf{Corrección de rumbo:} esto construye el dataset, así que va en
la sección de datos, no en la de modelamiento.

\subsection*{Modelamiento: las tres decisiones del comprador (slide 4)}
\begin{table}[h]
\centering
\begin{tabular}{p{4cm}p{3cm}p{7.5cm}}
\toprule
\textbf{Decisión} & \textbf{Qué es} & \textbf{Cómo se operacionaliza} \\
\midrule
``¿en qué cadena?'' & Modelo A, disponible hoy & Brecha de precio entre cadenas (regresión) o cadena más barata (multiclase). Features: categoría, marca, marca propia, gramaje, unidad, nivel de precio. Sin usar precios de otras cadenas. SHAP para interpretar. 11\,019 EAN. \\
``¿compro ahora o espero?'' & Modelo B, noviembre & $P(\texttt{on\_sale}$ en $t{+}1..t{+}7)$. Features: días desde la última oferta, precio vs.\ mediana histórica, si la cadena hermana ya bajó, día de semana, quincena. Requiere panel acumulado: es el modelo que justifica recolectar a diario. \\
``¿la quincena alcanza?'' & Canasta en soles & Analytics, entregable de cierre. \\
\bottomrule
\end{tabular}
\end{table}

Modelo C (opcional): duración de la oferta como análisis de supervivencia,
conecta con Bueno (2024).

Señal promocional disponible: 24.4\,\% de las filas en oferta, con descuento
mediano de 13--17\,\%; la intensidad varía 2.5$\times$ entre cadenas (Vivanda
13\,\%, Plaza Vea 32\,\%).

\section{Advertencias y pendientes}

\begin{itemize}[nosep]
  \item \textbf{No usar} la rigidez de 0.48\,\% ni la tasa de entrada a oferta de
        0.11\,\%: se midieron sábado$\to$domingo. Remedir lunes--viernes; decide
        si el Modelo B es viable a 7 días o hay que ampliar a 14--30.
  \item Corregir slides (Makro fuera, Vivanda dentro, 5 cadenas). Correr el
        backfill de EAN.
  \item La \textbf{regla de precios faltantes} (\S 2 del README) sigue abierta y es
        la decisión metodológica más importante.
  \item Confirmar si el curso exige deep learning o acepta ML clásico.
  \item Decidir si el Modelo C entra o queda como extensión.
  \item El 15\,\% de los pares con EAN discrepa en gramaje; los desacuerdos
        modelo-vs-EAN detectan errores de catalogación.
\end{itemize}

\subsection*{Horizonte}
\begin{tabular}{ll}
\toprule
Esta semana & Slides, backfill de EAN, tasa de cambio en días hábiles. \\
Octubre & Canasta sobre panel acumulado; regla de precios faltantes; Modelo A. \\
Noviembre & Rigidez y liderazgo de precios; Modelo B con 6+ semanas. \\
Diciembre & Canasta en soles + IPC-Web como anexo descriptivo, sin nowcast. \\
\bottomrule
\end{tabular}

\end{document}
